 \documentclass[12pt]{article}
 % LaTeX file created by GUIDE version 46.2 on 02/21/26 at 15:54
 %%%% Usage: latex tree.tex, dvips tree.tex, ps2pdf tree.tex
 %%%% or: xelatex tree.tex
 \usepackage{pstricks}
 %%%% If using pdflatex, replace above line with \usepackage[pdf]{pstricks}
 %%%% and use: pdflatex -shell-escape tree.tex
 \usepackage{pstricks,pst-node,pst-tree}
 \usepackage{geometry}
 \usepackage{xcolor}
 \usepackage{lscape}
 \definecolor{wheat}{rgb}{0.96, 0.87, 0.7}
 \definecolor{lightgray}{rgb}{0.9, 0.9, 0.9}
 \definecolor{aqua}{rgb}{0.0, 1.0, 1.0}
 \definecolor{mauve}{rgb}{0.88, 0.69, 1.0}
 \definecolor{skyblue}{RGB}{86,180,233}
 \definecolor{vermillion}{RGB}{213,94,0}
 \definecolor{purple}{RGB}{204,121,167}
 \definecolor{bluishgreen}{RGB}{0,158,115}
 \pagestyle{empty}
 \begin{document}
 %\begin{landscape}
 \begin{center}
\psset{linecolor=black,tnsep=1pt,tndepth=0cm,tnheight=0cm,treesep=1.5cm,levelsep=55pt,radius=10pt,fillstyle=solid}
  \pstree[treemode=D]{\Tcircle{ 1 }~[tnpos=l]{\shortstack[r]{\texttt{\detokenize{V1}}\\=\texttt{\detokenize{MUTATED}}}}
 }{
  \pstree[treemode=D]{\Tcircle{ 2 }~[tnpos=l]{\shortstack[r]{\texttt{\detokenize{V2}}\\=\texttt{\detokenize{POST_T}}}}
 }{
    \Tcircle[fillcolor=skyblue]{ 4 }~{\shortstack[c]{\emph{18}\\0}}  
    \Tcircle[fillcolor=yellow]{ 5 }~{\shortstack[c]{\emph{500}\\3.47}}  
   }
  \pstree[treemode=D]{\Tcircle{ 3 }~[tnpos=l]{\shortstack[r]{\texttt{\detokenize{V3}}\\=\texttt{\detokenize{POST_T}}}}
 }{
    \Tcircle[fillcolor=skyblue]{ 6 }~{\shortstack[c]{\emph{71}\\0}}  
    \Tcircle[fillcolor=skyblue]{ 7 }~{\shortstack[c]{\emph{420}\\1.86}}  
 }
 }
 \end{center}
GUIDE v.46.2 0.250-SE
piecewise-constant least-squares regression tree
for predicting \texttt{\detokenize{EXPRESSION_CLUSTER}}.
 % Tree constructed from 1009 observations
 % (excluding observations with non-positive weight or missing values
 % in \texttt{D}, \texttt{T}, \texttt{R} or \texttt{Z} variables)
 % and pruned from a tree with 318 terminal nodes.
 % Maximum number of split levels is 15 and minimum node sample size is 2.
At each split, an observation goes to the left branch 
 if and only if the condition is satisfied.
\texttt{\detokenize{V1}} = \texttt{\detokenize{IGHV_MUTATION_STATUS}}.
\texttt{\detokenize{V2}} = \texttt{\detokenize{TREATMENT_STATUS}}.
\texttt{\detokenize{V3}} = \texttt{\detokenize{TREATMENT_STATUS}}.
Sample size (in \emph{italics}) and mean of \texttt{\detokenize{EXPRESSION_CLUSTER}} printed below nodes.
 Terminal nodes with means above and below value of  2.49 at root node are painted \colorbox{yellow}{yellow} and \colorbox{skyblue}{skyblue} respectively.
 Second best split variable at root node is \texttt{\detokenize{IGHV_IDENTITY_PERCENTAGE}}.
 %\end{landscape}
 \end{document}
